\subsection{抛物子群}

定义 2. 若 G 的一个子群包含 B 的某个共轭子群，则称其为抛物子群。

显然，每个包含一个抛物子群的子群都是抛物子群。

命题 4. 设 P 是 G 的一个子群。

\begin{enumerate}
\def\labelenumi{\alph{enumi})}
\item
  P 是抛物子群，当且仅当存在 S 的一个子集 X，使得 P 与 \(G_{X}\) 共轭（关于 \(G_{X}\) 的定义见~第 5 号）。
\item
  设 \(\mathrm{X}, \mathrm{X}' \subset \mathrm{S}\) 且 \(g, g' \in \mathrm{G}\) 满足 \(\mathrm{P} = g\mathrm{G}_{\mathrm{X}}g^{-1} = g'\mathrm{G}_{\mathrm{X}'}g'^{-1}\)。则 \(\mathrm{X} = \mathrm{X}'\) 且 \(g'g^{-1} \in \mathrm{P}\)。
\end{enumerate}

断言 a) 由定理 3 b) 得出。

在 \(b\) 的假设下，有

\[
g ^ {- 1} g ^ {\prime} \mathrm{B} g ^ {\prime - 1} g \subset g ^ {- 1} g ^ {\prime} \mathrm{G} _ {\mathrm{X} ^ {\prime}} g ^ {\prime - 1} g = \mathrm{G} _ {\mathrm{X}},
\]

且命题 3 表明 \(g^{-1}g' \in \mathrm{G}_X\)。因此，\(\mathrm{G}_{\mathrm{X}'} = \mathrm{G}_X\) 且 \(\mathrm{X}' = \mathrm{X}\)，由定理 3 b)。最后，

\[
g ^ {\prime} g ^ {- 1} = g. g ^ {- 1} g ^ {\prime}. g ^ {- 1} \in g \mathrm{G} _ {\mathrm{X}} g ^ {- 1},
\]

这就证明了 b)。

若抛物子群 P 与 \(G_{X}\) 共轭，其中 \(X \subset S\)，则称 P 为 X 型。

定理 4. (i) 设 \(\mathbf{P}_1\) 和 \(\mathbf{P}_2\) 是 \(\mathbf{G}\) 的两个抛物子群，且其交为抛物子群；设 \(g \in \mathbf{G}\) 满足 \(g\mathrm{P}_1g^{-1} \subset \mathrm{P}_2\)。则 \(g \in \mathbf{P}_2\) 且 \(\mathbf{P}_1 \subset \mathbf{P}_2\)。

\begin{enumerate}
\def\labelenumi{(\roman{enumi})}
\setcounter{enumi}{1}
\item
  交为抛物子群的两个抛物子群不共轭。
\item
  设 \(\mathbf{Q}_1\) 和 \(\mathbf{Q}_2\) 是 \(\mathbf{G}\) 的两个抛物子群，包含于子群 \(\mathbf{Q}\) 中，而该群是 \(\mathbf{G}\) 的子群。则任意 \(g \in \mathbf{G}\)，若满足 \(g\mathbf{Q}_1g^{-1} = \mathbf{Q}_2\)，就属于 \(\mathbf{Q}\)。
\item
  每个抛物子群都是自身的规范化子 \(^{6}\)。
\end{enumerate}

断言 (i) 由命题 3 和 4 得出，并蕴含 (ii)。在 (iii) 的假设下，有 \(gQ_{1}g^{-1} \subset Q\)，由 (i) 得 \(g \in Q\)。最后，(iv) 由 (iii) 取 \(Q_{1} = Q_{2} = Q\) 得出。

命题 5. 设 \(\mathbf{P}_1\) 和 \(\mathbf{P}_2\) 是 \(\mathbf{G}\) 的两个抛物子群。则 \(\mathbf{P}_1 \cap \mathbf{P}_2\) 包含 \(\mathbf{T}\) 的某个共轭子群。

先用 G 的一个内自同构变换 \(P_{1}\) 和 \(P_{2}\)，可假定 \(B \subset P_{1}\)。设 \(g \in G\) 满足 \(gBg^{-1} \subset P_{2}\)。由定理 1，存在 \(n \in N\) 和 \(b, b' \in B\)，使得 \(g = bnb'\)。由于 T 在 N 中正规，

\[
\mathrm{P} _ {2} \supset g \mathrm{B} g ^ {- 1} = b n \mathrm{B} n ^ {- 1} b ^ {- 1} \supset b n \mathrm{T} n ^ {- 1} b ^ {- 1} = b \mathrm{T} b ^ {- 1}
\]

并且

\[
\mathrm{P} _ {1} \supset \mathrm{B} \supset b \mathrm{T} b ^ {- 1},
\]

这就证明了命题。

